\section{Прямой атомистический расчёт сольвуса}
\label{sec:ensemble}

\subsection{Степени свободы и механические условия}
В больших ячейках одновременно менялись химические типы узлов, координаты
атомов и объём. В основных сериях независимо менялось отношение осей
\(\eta=c/a\) при гексагональной симметрии. Для контроля геометрии C46
объём также релаксировал, а \(\eta\) фиксировалось. В обоих случаях
во всех вероятностях принятия использовано одно внешнее давление \(p=0\).
Две фазы могут иметь разные равновесные объёмы даже в отсутствие давления.

Чтобы область допустимых смещений двигалась вместе с кристаллом,
для каждого узла взяли вектор смещения \(\mathbf u_i\) относительно
идеального положения в текущей ячейке и вычли среднее
\(\overline{\mathbf u}\). Для двенадцати векторов \(\mathbf d_{ij}\)
от узла к его ближайшим соседям вычисляли
\[
 W=\max_{i,j}\frac{2(\mathbf u_i-\overline{\mathbf u})\cdot\mathbf d_{ij}}
                              {|\mathbf d_{ij}|^2}.
\]
Ход допустим при \(W<W_{\max}\). Условие сохраняли для обоих химических
типов и всех окон. Предложения объёма и формы дополнительно удовлетворяли
\(15<V/N<21\)~\AA$^3$/атом, \(1{,}45<\eta<1{,}72\) при свободных осях
и \(d_{\rm nn}(1-W_{\max})\ge 1{,}68052\)~\AA, где \(d_{\rm nn}\) —
кратчайшее расстояние между идеальными узлами текущей ячейки.
Так сохранялась область, для которой проведена проверка коротких расстояний.
В основной ячейке с \(W_{\max}=0{,}40\) у обеих фаз доля кадров в
пределах \(0{,}25\)~\AA$^3$/атом от границ объёма и \(0{,}01\) от
границ \(\eta\) была нулевой.

\subsection{Химический порядок и зонтичные окна}
Пусть \(n_j\) — число атомов Al на подрешётке \(j\), а \(k=\sum_{j=1}^4 n_j\).
Параметр порядка и энергия зонтика в окне \(l\) равны
\begin{equation}
 Q=\frac{4(4\max_j n_j-k)}{3N},\qquad
 B_l(Q)=\tfrac12 K_l(Q-q_l)^2 .
\end{equation}
У идеального D0$_{19}$ получаем \(Q=1\). В случайном растворе максимум
из четырёх заполнений даёт положительный \(Q\) при конечном \(N\).
Определение объединяет четыре трансляционно эквивалентных варианта
химического порядка. Упорядоченные домены разных
вариантов могут снижать общий \(Q\), поэтому проверяли также статистику
состава и распределение заполнений подрешёток.

Для химической замены в случайном узле вероятность принятия определяется
разностью \(U_{\rm MTP}+B_l-\Delta\mu_l k\), где \(\Delta\mu_l\) —
конфигурационное поле, положительное в пользу Al. Каждый проход включает
по \(N\) предложений замены и декартова смещения. Объём предлагается раз
за проход симметрично по \(V\). Для изотропного масштабирования координат
логарифм отношения весов содержит
\[
 -\beta[\Delta U_{\rm MTP}+p\,\Delta V]+N\ln(V'/V),
 \qquad \beta=(k_{\rm B}T)^{-1}.
\]
Относительные координаты после удаления переноса дают степень \(N-1\),
интегрирование общего переноса возвращает ещё множитель \(V\).
Для свободных осей каждые три прохода предлагали симметричное изменение
\(\ln\eta\) при постоянном объёме. В расчёте химических замен использованы
одинаковые фиктивные массы. Для физических классических масс
\(\Delta\mu_{\rm phys}=\Delta\mu_{\rm cfg}
 -\tfrac32 k_{\rm B}T\ln(m_{\rm Al}/m_{\rm Ti})\).
При 1000 К поправка составляет \(+0{,}07410\) эВ на замену.
Линейная по составу поправка к свободной энергии смещает наклон общей
касательной и сохраняет составы в точках касания.

\subsection{Как из окон получены две фазовые границы}
Центры зонтиков охватывали \(Q=0{,}035\)--\(0{,}910\).
Соседние окна перекрываются по заполнениям. При одинаковых \(N\),
\(W_{\max}\) и механических условиях общая атомистическая энергия и
мера координат сокращаются в отношении весов двух окон. Для их объединения
методом MBAR достаточно точно знать \(k,Q\) и добавленный член
\(\beta[B_l(Q)-\Delta\mu_l k]\) в каждом окне~\cite{mbar2008}.
На практике кадры с одним и тем же дискретным состоянием \(s=(k_s,Q_s)\)
объединяли, сохраняя общую кратность \(h_s\). Пусть \(M_l\) — число
производственных кадров окна \(l\), а \(\mu_0\) — медиана полей окон,
выбранная для удобства вычислений. Определим безразмерные добавки
\(u_l(s)=\beta[B_l(Q_s)-(\Delta\mu_l-\mu_0)k_s]\).
Относительные свободные энергии окон \(f_l\) определяются из уравнений MBAR
\begin{equation}
 D_s=\sum_l M_l\exp[f_l-u_l(s)],\qquad
 e^{-f_j}=\sum_s \frac{h_s}{D_s}\,e^{-u_j(s)},\qquad f_1=0 .
 \label{eq:umbrella-mbar}
\end{equation}
В программе решается эквивалентная задача минимизации с проверкой
градиента. Для целевого поля \(\Delta\mu'\) вес состояния после снятия
всех зонтиков равен
\[
 a_s(\Delta\mu')=\frac{h_s}{D_s}\exp[\beta(\Delta\mu'-\mu_0)k_s],
 \qquad P_s(\Delta\mu')=\frac{a_s(\Delta\mu')}{\sum_t a_t(\Delta\mu')}.
\]
Обе фазы получают общую нормировку через цепочку перекрывающихся окон.
Сглаживание применялось только при поиске пиков и минимума \(P(Q)\).
Для этого массу вероятности разбивали на интервалы шириной \(0{,}005\)
и сглаживали гауссовым ядром с шириной \(0{,}04\) по \(Q\).
Разделитель \(Q^\ddagger\) выбирали в минимуме между модами
при поле, где их высоты равны. С этим разделителем находили \(\Delta\mu^*\)
из условия
\[
 \frac{P_{\Delta\mu^*}(Q>Q^\ddagger)}
      {P_{\Delta\mu^*}(Q\le Q^\ddagger)}=4 .
\]
Четвёрка задаёт одинаковый термодинамический вес для неупорядоченной фазы
и каждого из четырёх вариантов D0$_{19}$~\cite{janke1993}.
При найденном поле левый и правый условные средние \(\langle k/N\rangle\)
дают составы \(\alpha\) и \(\alpha_2\). Пересечение изотерм
\(\langle x\rangle(\Delta\mu)\) разных размеров служит независимой по
критерию проверкой положения перехода~\cite{borgs1992}.

Для \(N=1400,\ W<0{,}40\) и свободного \(c/a\) использованы по 14 окон
при полях \(3{,}287\) и \(3{,}290\) эВ на замену. Жёсткость каждого
зонтика \(K_l=94{,}6\) эВ относится ко всей ячейке. Получено 22\,400
кадров. Равенство высот найдено при \(3{,}29258\) эВ, минимум между
пиками — при \(Q^\ddagger=0{,}4975\). При этом разделителе равенство
весов на один вариант достигается при \(3{,}29368\) эВ. Эффективное
число взвешенных кадров Киша, \(1/\sum_i w_i^2\), для нормированных
внутри каждой фазы покадровых весов \(w_i\) равно примерно 3510 и 3325.
Поправка на корреляцию последовательных кадров даёт ориентировочно
51 и 64 соответственно. Поэтому ошибка средних по фазам оценивалась
по временным блокам.

Статистические интервалы рассчитаны ресэмплированием последовательных
неперекрывающихся блоков внутри каждой траектории. Целевую длину блока
брали равной удвоенному максимуму статистической неэффективности
для \(k,Q,V,c/a\), округляя вверх до целого числа кадров.
В указанной серии это дало от 3 до 22 блоков на траекторию. В каждой бутстрэп-выборке
заново оценивали веса, разделитель и поле. Для \(N=96,288,768,1400\)
использованы соответственно 96, 128, 192 и 192 бутстрэп-реализации.
Эти интервалы характеризуют шум данной выборки. Зависимость от размера,
чувствительность к \(W_{\max}\) и разброс независимых стартов
приведены отдельно.
